\documentclass[../../../include/open-logic-section]{subfiles}

\begin{document}

\olfileid{sth}{replacement}{refproofs}
\olsection{ملحق: نتائج تتعلق بالاستبدال}

نثبت في هذا القسم الانعكاس داخل~$\ZF$، ثم نبين وجهًا يصير به الانعكاس
مكافئًا للاستبدال، ونختم بنتيجة لافتة في قوتهما. \emph{تنبيه: هذا
القسم، بفارق ظاهر، أرفع أجزاء الكتاب منزلة في الصعوبة الرياضية.}

ونقدم لمّة نستعمل فيها، اختصارًا، وسيلة ترميزية هي \emph{الخط
العلوي} للدلالة على متتاليات المتغيرات أو الأشياء. فالعبارة
«$\overline{a}_k$» تغني عن «$a_{k_1}$، \dots، $a_{k_n}$»، ويعين
السياق قيمة~$n$.

\begin{lem}\ollabel{lemreflection}
لكل~$1 \leq i \leq k$، لتكن~$\phi_i(\overline{v}_i, x)$ صيغة. عندئذ
يوجد لكل~$\alpha$ عدد~$\beta > \alpha$ بحيث، لأي
$\overline{a}_1, \ldots, \overline{a}_k \in V_\beta$ ولكل
$1 \leq i \leq k$:
\[
	\exists x\phi_i(\overline{a}_i, x) \rightarrow (\exists x \in V_\beta) \phi_i(\overline{a}_i, x)
\]
\end{lem}

\begin{proof}
نعرّف حدًا~$\mu$ كما يلي: نجعل
$\mu(\overline{a}_1, \ldots, \overline{a}_k)$ أول مرحلة~$V$ تحقق،
لكل~$1 \leq i \leq k$، الشرطيات الآتية:
\[
\exists x\phi_i(\overline{a}_i, x) \rightarrow (\exists x \in V) \phi_i(\overline{a}_i, x)
\]
ولا يعسر التحقق من وجود~$\mu(\overline{a}_1, \ldots,
\overline{a}_k)$ لكل~$\overline{a}_1, \ldots, \overline{a}_k$.
ثم نستخدم الاستبدال ومبرهنة العودية فنعرّف:
\begin{align*}
	S_0 & = V_{\alpha+1}\\
	S_{n+1} & = S_n \cup \bigcup
	\Setabs{\mu(\overline{a}_1, \ldots, \overline{a}_k)}
	{\overline{a}_1, \ldots, \overline{a}_k \in S_n} \\
	S &= \bigcup_{m < \omega} S_m.
\end{align*}
كل~$S_n$، ومن ثم~$S$ نفسها، مرحلة تلي~$V_\alpha$. والآن ثبت
$\overline{a}_1$، \dots،~$\overline{a}_k \in S$؛ فلا بد من
$n < \omega$ بحيث~$\overline{a}_1$، \dots،
$\overline{a}_k \in S_n$. وثبت~$1 \leq i \leq k$، وافترض
$\exists x\phi_i(\overline{a}_i,x)$. يقتضي البناء
$(\exists x \in \mu(\overline{a}_1,
\ldots, \overline{a}_k))\phi_i(\overline{a}_i, x)$؛ ومنه
$(\exists x \in S_{n+1})\phi_i(\overline{a}_i, x)$، ثم
$(\exists x \in S)\phi_i(\overline{a}_i, x)$. فهذه~$S$ هي
$V_\beta$ المطلوبة.
\end{proof}
\noindent
وبهذه اللمّة يصير إثبات
\olref[sth][replacement][ref]{reflectionschema} مباشرًا:

\begin{proof}[الإثبات]
ثبّت~$\alpha$. ولا يخل بالعموم أن نجعل روابط~$\phi$ مقصورة على
$\exists$ و$\lnot$ و$\land$، فهي وافية من جهة القوة التعبيرية.
ولتكن~$\psi_1, \ldots, \psi_k$ جميع الصيغ الجزئية لـ$\phi$، مرتبة
بحسب تعقيدها، بحيث~$\psi_k = \phi$. ومن~\olref{lemreflection} يوجد
$\beta > \alpha$ بحيث، لأي~$\overline{a}_i \in V_\beta$ ولكل
$1 \leq i \leq k$:
\begin{align}\label{reflectionnicelybehaved}
	\exists x\psi_i(\overline{a}_i, x) \rightarrow
	(\exists x \in V_\beta) \psi_i(\overline{a}_i, x)\tag{*}
\end{align}
ونثبت، بالاستقراء في تعقيد~$\psi_i$، أنه لأي
$\overline{a}_i \in V_\beta$:
$\psi_i(\overline{a}_i) \leftrightarrow
\psi_i^{V_\beta}(\overline{a}_i)$. أما إذا كانت~$\psi_i$ ذرية
فالأمر ظاهر. ويثبت هذا التكافؤ كذلك أنه إذا كانت~$\psi_i$ نفيًا أو
اقترانًا لصيغ جزئية تحققه، فإن~$\psi_i$ نفسها تحققه؛ فلا تبقى إلا
حالة التكميم. ثبت
$\overline{a}_i \in V_\beta$؛ فعندئذ:
\begin{align*}
	(\exists x \psi_i(\overline{a}_i, x))^{V_\beta}
	&\text{ إذا وفقط إذا }
	(\exists x \in V_\beta)\psi_i^{V_\beta}(\overline{a}_i, x)
	&&\text{بحكم التعريف}\\
	&\text{ إذا وفقط إذا }
	(\exists x \in V_\beta)\psi_i(\overline{a}_i,  x)
	&&\text{بحكم فرض الاستقراء}\\
	&\text{ إذا وفقط إذا }
	\exists x \psi_i(\overline{a}_i, x)
	&&\text{بموجب \eqref{reflectionnicelybehaved}}
\end{align*}
فتم الاستقراء، والنتيجة تابعة من~$\psi_k = \phi$.
\end{proof}

ثبت إذن الانعكاس في~$\ZF$، وبرهاننا في جوهره طريق
\citet{Montague1961}. ونريد الآن أن نثبت، داخل~$\Z$، أن الانعكاس
يستلزم الاستبدال. والبرهان الآتي تبسيط لطريق~\citet{Levy1960}.

ولأننا نعمل في~$\Z$، لا يمكننا عرض الانعكاس بالصورة السابقة؛ فقد
استعملت تلك الصورة الرمز «$V_\alpha$»، وهو غير قابل للتعريف في~$\Z$
(انظر~\olref[sth][spine][zf]{sec}). فنستعيض عنها بصورة تبدو أضعف:

\begin{defish}
\emph{الانعكاس الضعيف.} لكل صيغة~$\phi$، توجد مجموعة متعدية~$S$ بحيث
تكون~$0$ و$1$ وأي وسائط لـ$\phi$ من !!{element}s~$S$، ويكون
$(\forall \overline{x} \in S)(\phi \liff \phi^S)$.
\end{defish}

ولكي نستعمل هذا في إثبات الاستبدال، نتبع أولًا
\citet[first part of Theorem 2]{Levy1960} ونثبت إمكان «عكس» صيغتين
معًا:

\begin{lem}[في $\Z + \text{الانعكاس الضعيف}$.]\ollabel{lem:reflect}
لكل صيغتين~$\psi, \chi$، توجد مجموعة متعدية~$S$ بحيث يكون~$0$ و$1$
(وأي وسائط للصيغتين) من !!{element}s~$S$، ويكون
$(\forall \overline{x} \in S)((\psi \liff \psi^S) \land (\chi
\liff \chi^S))$.
\end{lem}

\begin{proof}
ضع~$\phi$ مساوية للصيغة~$(z = 0 \land \psi) \lor (z = 1 \land \chi)$.

وفي هذا اختصار ينبغي كشفه: فالعبارة «$z = 0$» حقها أن تفصل إلى
«$\forall t\, t \notin z$»، والعبارة «$z =1$» إلى «$\forall s(s \in z
\liff \forall t\, t \notin s)$». غير أن~$0, 1 \in S$، وأن~$S$
متعدية؛ فهاتان الصيغتان \emph{مطلقتان} بالنسبة إلى~$S$، أي تصدقان في
الشيء نفسه، سواء قيدنا مكمماتهما بـ$S$ أم أطلقناها.
\footnote{وبصيغة أدق، إذا كانت~$\xi$ إحدى هاتين الصيغتين، فإن
$\xi(z) \liff \xi^S(z)$.}

ويعطينا الانعكاس الضعيف~$S$ ملائمة بحيث:
\begin{align*}
	(\forall z, \overline{x} \in S)(&\phi \liff \phi^S)\\
	\text{أي }(\forall z, \overline{x} \in S)(&((z = 0 \land \psi) \lor (z = 1 \land \chi)) \liff {}\\
	&\phantom{(}((z = 0 \land \psi) \lor (z = 1 \land \chi))^S)\\
	\text{أي }(\forall z, \overline{x} \in S)(&((z = 0 \land \psi) \lor (z = 1 \land \chi))\liff {}\\
	&\phantom{(}((z = 0 \land \psi^S) \lor (z = 1 \land \chi^S)))\\
	\text{أي }(\forall \overline{x} \in S)(&(\psi \liff \psi^S) \land (\chi \liff \chi^S))
\end{align*}
فالادعاء الثاني يقتضي الثالث لإطلاق «$z = 0$» و«$z=1$» بالنسبة
إلى~$S$؛ والرابع تابع لأن~$0 \neq 1$.
\end{proof}\noindent
ثم نحصل على الاستبدال باتباع
\citet[Theorem 6]{Levy1960} مع تبسيط الحجة:

\begin{thm}[في $\Z$ + الانعكاس الضعيف]\label{thm:replacement}
لكل صيغة~$\phi(v,w)$ ولكل~$A$، إذا كان
$(\forall x \in A)\lexists![y][\phi(x,y)]$، فإن
$\Setabs{y}{(\exists x \in A)\phi(x,y)}$ موجودة.
\end{thm}

\begin{proof}
ثبّت~$A$ بحيث~$(\forall x \in A)\lexists![y][\phi(x,y)]$، وعرّف
الصيغتين:
\begin{align*}
	\psi &\text{ هي } (\phi(x, z) \land A = A)\\
	\chi &\text{ هي } \lexists[y][\phi(x, y)]
\end{align*}
من~\olref{lem:reflect}، ولأن~$A$ وسيط في~$\psi$، توجد مجموعة متعدية
$S$ بحيث~$0, 1, A \in S$، ومعها سائر الوسائط، وبحيث:
\[
	(\forall x,z \in S)((\psi \liff \psi^S) \land (\chi \liff \chi^S))
\]
فيكون، على الخصوص:
\begin{align*}
	(\forall  x, z \in S)(&\phi(x, z) \liff \phi^S(x, z))\\
	(\forall x \in S)(&\exists y\phi(x, y) \liff (\exists y \in S)\phi^S(x, y))
\end{align*}
نجمع بينهما، ونلاحظ أن~$A \subseteq S$ لأن~$A \in S$ و$S$ متعدية،
فنحصل على:
\begin{align*}
	(\forall x \in A)(&\exists y\phi(x, y) \liff (\exists y \in S)\phi(x, y))
\end{align*}
ولما كان~$(\forall x \in A)\lexists![y \in S][\phi(x,y)]$، من
الفرض~$(\forall x \in A)\lexists![y][\phi(x, y)]$، أعطانا الفصل
$\Setabs{y \in S}{(\exists x \in A) \phi(x, y)} = \Setabs{y}{(\exists
x \in A) \phi(x, y)}$.
\end{proof}

\end{document}
